\section{Оценка эффективности}

\subsection{Описание эксперимента}

Эксперимент проведён в изолированной среде на базе CALDERA 4.0 и Proxmox VE 8.1, включающей:
\begin{itemize}
    \item 50 сценариев атак (Log4Shell, phishing, PowerShell lateral movement, DNS tunneling);
    \item 1000 событий легитимной активности;
    \item контрольный модуль на основе SIEM-правил (Sigma + Elastic Detection Rules).
\end{itemize}

\subsection{Результаты}

Метод на основе атрибутивных метаграфов показал:
\begin{itemize}
    \item Точность — 94,3\,\%, полнота — 89,7\,\%, ложные срабатывания — 3,1\,\%;
    \item Среднее время реакции — 1.8 с (против 5.4 с у SIEM);
    \item Устойчивость к увеличению фона: снижение F1-меры всего на 1.2\,\% при 10× шуме.
\end{itemize}

Сравнение с графами знаний \cite{li2022kg} показало преимущество по полноте (+11.5\,\%) и скорости анализа (на 220 мс быстрее).